\sectionkicker{Source-backed technical notes}
\subsection*{Paper Watch — agent評価・測定bias・unified multimodal: Technical Notes}
本欄は記事本文で圧縮した一次資料上の情報を、比較・再検証しやすい形で整理したものである。時系列、確認済みの主張、留意点、一次資料URLを掲載する。
\medskip
\subsection*{Theme at a glance}
\addcontentsline{toc}{subsection}{Theme at a glance}
\begin{center}
\footnotesize
\begin{tabularx}{\linewidth}{@{}p{0.20\linewidth}p{0.10\linewidth}p{0.14\linewidth}X@{}}
\toprule
資料 & 位置づけ & 種別 & 時系列 \\
\midrule
AgenticVBench: Can AI Agents Complete Real-World Post-Production Tasks? & 主要資料 & 論文 & 2026-05-26 (論文\_RELEASE) \\
Identifying and Mitigating Systemic Measurement Bias in Production LLM Inference Benchmarks & 主要資料 & 論文 & 2026-05-22 (論文\_RELEASE) \\
Representation Forcing for Bottleneck-Free Unified Multimodal Models & 主要資料 & 論文 & 2026-05-29 (論文\_RELEASE) \\
\bottomrule
\end{tabularx}
\end{center}
\normalsize

\begin{technicalnote}{AgenticVBench: Can AI Agents Complete Real-World Post-Production Tasks?}{主要資料}
\begin{tabularx}{\linewidth}{@{}>{\bfseries}p{0.22\linewidth}X@{}}
組織 & - \\
種別 & 論文 \\
時系列 & 2026-05-26 (論文\_RELEASE) \\
\end{tabularx}
\smallskip
{\bfseries 一次資料から整理したtechnical points}
\begin{itemize}[leftmargin=1.5em,itemsep=0.35em]
\item \textbf{一次情報で確認できる事実}: The preserved primary source identifies “AgenticVBench: Can AI Agents Complete Real-World Post-Production Tasks?” on 2026-05-26.
\item \textbf{Author claim}: AgenticVBench evaluates multimodal agents on real post-production workflows requiring text, image, audio, video, planning, and tool use.
\end{itemize}
{\bfseries 読む際の境界}
\begin{itemize}[leftmargin=1.5em,itemsep=0.35em]
\item \textbf{分析上の留意点}: The preserved original arXiv metadata/abstract is primary-paper evidence, but experimental results and generalization are author-reported and were not independently reproduced in this Evidence review.
\end{itemize}
{\bfseries 一次資料}
\begingroup\sloppy
\begin{itemize}[leftmargin=1.5em,itemsep=0.25em]
\item {\scriptsize\url{http://arxiv.org/abs/2605.27705v1}}
\end{itemize}
\endgroup
\end{technicalnote}
\medskip

\begin{technicalnote}{Identifying and Mitigating Systemic Measurement Bias in Production LLM Inference Benchmarks}{主要資料}
\begin{tabularx}{\linewidth}{@{}>{\bfseries}p{0.22\linewidth}X@{}}
組織 & - \\
種別 & 論文 \\
時系列 & 2026-05-22 (論文\_RELEASE) \\
\end{tabularx}
\smallskip
{\bfseries 一次資料から整理したtechnical points}
\begin{itemize}[leftmargin=1.5em,itemsep=0.35em]
\item \textbf{一次情報で確認できる事実}: The preserved primary source identifies “Identifying and Mitigating Systemic Measurement Bias in Production LLM Inference Benchmarks” on 2026-05-22.
\item \textbf{Author claim}: The paper argues that production LLM inference benchmarks can acquire systemic client-side measurement bias and models Python GIL effects on TTFT/TPOT using queueing analysis.
\end{itemize}
{\bfseries 読む際の境界}
\begin{itemize}[leftmargin=1.5em,itemsep=0.35em]
\item \textbf{分析上の留意点}: The preserved original arXiv metadata/abstract is primary-paper evidence, but experimental results and generalization are author-reported and were not independently reproduced in this Evidence review.
\end{itemize}
{\bfseries 一次資料}
\begingroup\sloppy
\begin{itemize}[leftmargin=1.5em,itemsep=0.25em]
\item {\scriptsize\url{http://arxiv.org/abs/2605.24217v2}}
\end{itemize}
\endgroup
\end{technicalnote}
\medskip

\begin{technicalnote}{Representation Forcing for Bottleneck-Free Unified Multimodal Models}{主要資料}
\begin{tabularx}{\linewidth}{@{}>{\bfseries}p{0.22\linewidth}X@{}}
組織 & - \\
種別 & 論文 \\
時系列 & 2026-05-29 (論文\_RELEASE) \\
\end{tabularx}
\smallskip
{\bfseries 一次資料から整理したtechnical points}
\begin{itemize}[leftmargin=1.5em,itemsep=0.35em]
\item \textbf{一次情報で確認できる事実}: The preserved primary source identifies “Representation Forcing for Bottleneck-Free Unified Multimodal Models” on 2026-05-29.
\item \textbf{Author claim}: Representation Forcing makes visual representations explicit autoregressive intermediate targets inside a unified multimodal model, avoiding a separately pretrained generative VAE.
\end{itemize}
{\bfseries 読む際の境界}
\begin{itemize}[leftmargin=1.5em,itemsep=0.35em]
\item \textbf{分析上の留意点}: The preserved original arXiv metadata/abstract is primary-paper evidence, but experimental results and generalization are author-reported and were not independently reproduced in this Evidence review.
\end{itemize}
{\bfseries 一次資料}
\begingroup\sloppy
\begin{itemize}[leftmargin=1.5em,itemsep=0.25em]
\item {\scriptsize\url{http://arxiv.org/abs/2605.31604v4}}
\end{itemize}
\endgroup
\end{technicalnote}
\medskip
